% ============================================================
%  5. OBJETIVOS   (guía, 2.5)
%  Coherencia exigida: CADA objetivo específico debe tener
%  evidencia de cumplimiento dentro del desarrollo (sección 7)
%  y debe contribuir al objetivo general.
% ============================================================
\chapter{Objetivos}
\label{cap:objetivos}

\section{Objetivo General}
\label{sec:objetivo-general}

Construir un diagnóstico espacial de las brechas de información de transporte público en el
eje metropolitano de Cochabamba, a partir de las consultas de ruta de Trufi App del período
septiembre de 2022 a junio de 2024, y derivar de él una lista priorizada de zonas para
orientar el trabajo de mapeo de la organización.

\section{Objetivos Específicos}
\label{sec:objetivos-especificos}

Los objetivos específicos descomponen el objetivo general en resultados parciales y
verificables, alineados con el proceso de trabajo:

\begin{enumerate}
  \item [OE1.] Consolidar y auditar los 85 archivos semanales de consultas en un conjunto
        único y trazable, documentando esquema, calidad, duplicados, vacíos temporales y
        procedencia de las variables derivadas recibidas. \textbf{(Completado --- ver
        sección~\ref{sec:comprension-datos}.)}
  \item [OE2.] Caracterizar la demanda de información mediante análisis exploratorio
        espacio-temporal sobre una teselación H3, evaluando la sensibilidad de los
        resultados a la resolución espacial elegida.
  \item [OE3.] Construir, por celda territorial y semana, indicadores que integren demanda
        de consultas, cobertura de rutas GTFS, centralidad respecto del núcleo urbano y
        composición temporal del uso, documentando el método de cálculo de cada uno y las
        variables descartadas por fuga de información.
  \item [OE4.] Contrastar formalmente, mediante pruebas de hipótesis, si la brecha de
        cobertura de transporte es mayor en la periferia que en el centro, y si un modelo
        de predicción de demanda mejora de forma significativa sobre líneas base
        explícitas.
  \item [OE5.] Entrenar y comparar modelos de predicción de demanda de consultas por celda
        y período frente a líneas base explícitas, con validación cronológica.
  \item [OE6.] Derivar y validar una regla de priorización de mapeo que combine demanda
        esperada y brecha de cobertura, evaluando con un criterio declarado de antemano si
        el modelo de predicción aporta valor de ordenamiento frente a una línea base simple.
  \item [OE7.] Proponer una arquitectura de actualización periódica del diagnóstico, con un
        prototipo ejecutable y un plan de monitoreo, para su uso operativo por parte de
        Trufi App.
\end{enumerate}

\vspace{6pt}
\noindent
Nota de coherencia: cada objetivo específico se corresponde con una fase del desarrollo
presentado en el capítulo~\ref{cap:desarrollo} y es retomado explícitamente en las
conclusiones del capítulo~\ref{cap:conclusiones}. La tabla~\ref{tab:matriz-coherencia}
resume esa correspondencia y el estado de avance de cada objetivo.

\begin{table}[htbp]
  \centering
  \caption{Matriz de coherencia objetivos-resultados-evidencias}
  \label{tab:matriz-coherencia}
  \interlineadoSimple
  \begin{tabularx}{\textwidth}{@{}p{0.28\textwidth}p{0.32\textwidth}X@{}}
    \toprule
    \textbf{Objetivo específico} & \textbf{Resultado esperado} & \textbf{Evidencia en la monografía} \\
    \midrule
    OE1. Auditar y consolidar datos & Conjunto de datos único, trazable y con calidad documentada & Sección~\ref{sec:comprension-datos} (completa) \\
    OE2. Caracterizar demanda (EDA espacial) & Tablas y figuras de distribución espacio-temporal; análisis de sensibilidad de resolución H3 & Secciones~\ref{sec:comprension-datos} y~\ref{sec:preparacion-datos} (completa) \\
    OE3. Construir indicadores por celda-semana & Tabla de indicadores con demanda, cobertura GTFS, centralidad y composición temporal & Sección~\ref{sec:preparacion-datos} (completa) \\
    OE4. Contrastar hipótesis de brecha y de mejora predictiva & Pruebas de Mann-Whitney (H1) y Diebold-Mariano/Wilcoxon (H2) & Sección~\ref{sec:evaluacion-resultados} (completa) \\
    OE5. Predecir demanda & Cuatro modelos comparados contra dos líneas base, con métricas de error & Secciones~\ref{sec:modelado} y~\ref{sec:evaluacion-resultados} (completa) \\
    OE6. Regla de priorización de mapeo & Ranking de celdas por demanda no resuelta, con contraste preinscrito del aporte del modelo & Secciones~\ref{sec:evaluacion-resultados} y~\ref{sec:despliegue} (completa) \\
    OE7. Proponer despliegue & Prototipo ejecutable, arquitectura y plan de monitoreo & Sección~\ref{sec:despliegue} (completa) \\
    \bottomrule
  \end{tabularx}
  \fuente{Elaboración propia.}
\end{table}
